\subsection{広義積分の定義}

\sectionref{節:ベータ積分}では,端点を避けた積分の極限によって,ベータ積分を定めた.ここでは,無限区間や区間の内部に特異点がある場合にも積分の意味を広げる.

$\int_0^1\frac1{\sqrt{x}}\dd x$では,\,$x=0$で被積分関数が定義されず,その近くで値が限りなく大きくなる.それでも,\,$\varepsilon>0$で切り取った区間では
\[
 \int_{\varepsilon}^1\frac1{\sqrt{x}}\dd x
 =2-2\sqrt{\varepsilon}
\]
と通常の定積分として計算できる.同じことを$\frac1x$でも行うと$\int_{\varepsilon}^1\frac1x\dd x=-\log\varepsilon$である.\,$\varepsilon$を小さくしていくと次のようになる.
\begin{center}
\begin{tabular}{ccc}
\toprule
$\varepsilon$ & $\int_{\varepsilon}^1\frac1{\sqrt{x}}\dd x$ & $\int_{\varepsilon}^1\frac1x\dd x$\\
\midrule
$10^{-2}$ & $1.8$ & $4.605\cdots$\\
$10^{-4}$ & $1.98$ & $9.210\cdots$\\
$10^{-6}$ & $1.998$ & $13.815\cdots$\\
\bottomrule
\end{tabular}
\end{center}
どちらの関数も$x=0$の近くで限りなく大きくなるが,前者の積分は$2$に近づき,後者の積分は限りなく大きくなる.前者のように極限が有限の値になるとき,その値を積分の値と定めよう.

\begin{definition}{広義積分}
$f$が$[a,b)$の各閉部分区間で積分可能で,端点$b$に特異性をもつとき,
\[
 \int_a^bf(x)\dd x=\lim_{c\to b-}\int_a^cf(x)\dd x
\]
と定める.左端に特異性がある場合も同様に,一側極限で定める.

上端が無限大の場合は
\[
 \int_a^{\infty}f(x)\dd x=\lim_{R\to\infty}\int_a^Rf(x)\dd x
\]
と定める.下端が$-\infty$の場合も同様である.

いずれの場合も,右辺の極限が有限の実数として存在するとき,この広義積分は収束するといい,有限の極限が存在しないとき発散するという.
\end{definition}

ここで,\,$+\infty$へ発散する場合も,振動して極限がない場合も,「収束しない」という点では同じである.発散の仕方を区別したいときには「$+\infty$へ発散する」などと具体的に述べる.上の例では,\,$\int_0^1\frac1{\sqrt{x}}\dd x=2$と収束し,\,$\int_0^1\frac1x\dd x$は$+\infty$へ発散するのだ.
